%%%%%%%%%%  scan idx 60  |  folio 47  %%%%%%%%%%
\begin{equation}\tag{2.14}
 f(E)=g(\cO_X(E))
\end{equation}
\emph{pour tout fibré vectoriel $E$ sur $X$, en notant $\cO_X(E)$ le faisceau des germes de sections régulières de $E$ sur $X$.}

Pour calculer $g(F)$ pour un faisceau algébrique cohérent quelconque, on considère une résolution finie de $F$ par des faisceaux de la forme $\cO_X(E_i)$, et l'on a alors~:
\begin{equation}\tag{2.14 bis}
 g(F)=\sum_{i\geq0}(-1)^if(E_i).
\end{equation}

En particulier, si $X$ est connexe, on a envisagé au \S 1 une fonction additive $\Ctil(E)$ du fibré vectoriel $E$, à valeurs dans le groupe $\Atil=\Atil(X)$. On en déduit une fonction additive de faisceaux cohérents, soit $\Ctil(F)$, et en passant aux composantes, des fonctions $F\longrightarrow C^i(F)\in A^i(X)$. Les $C^i(F)$ sont appelées les \emph{classes de Chern du faisceau cohérent} $F$~; elles se calculent théoriquement à partir des classes de Chern des fibrés vectoriels par (2.14 bis) (mais où il faut passer au deuxième membre à l'écriture multiplicative). Quant à la composante de $\Ctil(F)$ suivant $\Z$, c'est le \emph{rang} du faisceau $F$, qui correspond au rang d'un module sur un anneau intègre~: si $X$ est affine (cas auquel on peut toujours se ramener), $F$ est défini par un module de type fini sur l'anneau de coordonnées de $X$, et l'on prend le rang de ce module\dots.

\subsection*{\S 3. Généralités fonctorielles sur $\underline{K}(X)$}

Si $X$ est un espace algébrique, on pose pour abréger $\underline{K}(X)=\underline{K}(\cF(X))$, et lorsque $X$ est quasi-projectif non singulier, on identifie ce groupe à $\underline{K}(\xi(X))$ en vertu du corollaire 2 au théorème 2.3. Comme $\underline{K}(\xi(X))$ est 
